\documentclass[10pt,a4paper]{amsart}
\usepackage[margin=19mm]{geometry}
\usepackage{amsmath,amssymb,mathtools}
\usepackage[scr=rsfs,cal=euler]{mathalfa}
\usepackage{xfrac}
\usepackage[unicode,hidelinks,bookmarks=false]{hyperref}
\newcommand{\ie}{i.e.\ }
\newcommand{\cf}{cf.\ }
\newcommand{\resp}{respectively\ }
\newcommand{\Subsection}{\subsection}
\providecommand{\nobreakdash}{\nobreak\hbox{-}}
\setlength{\emergencystretch}{2em}
\multlinegap0pt
\usepackage{cancel}
\usepackage{amssymb}
\usepackage{tikz}
\newtheorem{lemma}{Lemma}
\title{A proof-theoretic approach to abstract interpretation}
\author{Vijay D'Silva, Alessandra Palmigiano, Apostolos Tzimoulis, Caterina Urban}
\date{Source: arXiv:2605.26591v1 (CC BY 4.0)}
\begin{document}
\maketitle

\noindent\emph{Source and licence.} A proof-theoretic approach to abstract interpretation. Version arXiv:2605.26591v1 (CC BY 4.0), distributed by arXiv under the Creative Commons Attribution 4.0 International licence, http://creativecommons.org/licenses/by/4.0/, as recorded in the arXiv licence metadata for this exact version and shown on the abstract page https://arxiv.org/abs/2605.26591v1. Full licence notice: \texttt{licenses/arxiv-2605.26591-CC-BY-4.0.txt}.\par\smallskip

\noindent\textit{Editorial scope.} Counted unit: the article's lemma lemma (source0.tex lines 312--314). The article's complete original proof of it is delivered with it (source0.tex lines 315--321, one \texttt{proof} environment of 7 lines). No mathematics is written, repaired or re-derived: the statement and the proof are the article's own lines, in the article's own notation, and no retained line is substituted or re-flowed.  No further source-local context is needed: every cross-reference of every retained line resolves inside the two delivered spans.  Nothing was omitted from any retained span, so the package carries no omission record.  No retained line contains a citation, so the wrapper carries no bibliography and no bibliography entry.  No retained line contains a figure environment, an image-inclusion command or drawing code, so no image is delivered and the package declares no asset dependency.  Editorial formatting only, in addition: an AMS \texttt{amsart} preamble that restates the article's own declarations and macros used by the retained lines, namely the environments \texttt{lemma} on separate counters, together with the standard packages those lines need.  The retained environments print in this document as Lemma 1 in order of appearance, whereas the article prints the same items with the numbering of its own full document.  The complete native source of the article as distributed in the arXiv e-print for this exact version is bundled unchanged as \texttt{sources/201428/source0.tex}.
\begin{lemma}
	The algebra $\mathbb{L}$ is isomorphic to $\mathcal{A}$.
\end{lemma}

\begin{proof}
	Let $e:\mathcal{A}\to\mathbb{L}$ be defined as $e(a)=[a(x)]$, where $[a(x)]$ denotes the equivalence class of the predicate $a(x)$. First notice that $e$ is surjective since each element in the Lindenbaum-Tarski algebra of the logic $\mathcal{L}_A$ is  the equivalence class of some predicate. This follows immediately by induction on the complexity of the formulas and items 4 and 5 in the construction.
	
	To show that $e$ is injective let $[a(x)]=[b(x)]$. Then we have that $a(x)\dashv\vdash b(x)$. Notice that $\mathcal{L}_A$ is sound w.r.t.\ $\mathcal{A}$ when each predicate $c(x)$ is interpreted as $c$. Hence, we immediately get that $a=b$.
	
	Finally, $e$ being an order embedding and a homomorphism w.r.t.\ the connectives of the language follows immediately from items 4, 5 and 6, and the soundness of the logic w.r.t.\ $\mathcal{A}$.
\end{proof}

\end{document}
